% GENERATED FILE. Edit canonical lessons, then run npm run generate:books.
% canonical-source-hash: fnv1a64:2f89dfe6a38b54ad
% canonical-lessons: SA-C535-pratipalayati, SA-C535-satyasandhah, SA-C535-visvastah, SA-C535-bhangah, SA-C535-vacanam

\chapter{Keeping One's Word}
\label{ch:sa-keeping-one-s-word}
\hlchaptermodality{535}

\begin{chapteropening}
\textbf{By the end of this chapter:} \emph{I can say what happens if someone keeps or breaks their word.}

Complete the last lesson of chapter 535: I can say what happens if someone keeps or breaks their word.
\end{chapteropening}

\section[\texorpdfstring{\sk{प्रतिपालयति} (pratipālayati)}{प्रतिपालयति (pratipālayati)}]{\sk{प्रतिपालयति} (pratipālayati) — keeps to; abides by}

\label{lesson:SA-C535-pratipalayati}



\begin{quote}
Before the new one: say the Sanskrit for turns down, then the Sanskrit for consent.
\end{quote}

\subsection*{You'll want to know: \sk{प्रतिपालयति}}
\textbf{\sk{प्रतिपालयति}} — \emph{pratipālayati} — \textquotedblleft{}keeps to; abides by\textquotedblright{}.

\begin{quote}
\sk{यदि} \sk{सः} \sk{नियमं} \sk{प्रतिपालयति}, \sk{तर्हि} \sk{सर्वे} \sk{तं} \sk{विश्वसन्ति।}
\end{quote}

\begin{quote}
\emph{yadi saḥ niyamaṃ pratipālayati, tarhi sarve taṃ viśvasanti.}
\end{quote}

\begin{quote}
If he keeps the rule, everyone trusts him.
\end{quote}

Past (he …): \sk{सः} \sk{नियमं} \sk{प्रतिपालितवान्}, \emph{saḥ niyamaṃ pratipālitavān}.

\begin{cousinweb}
\textbf{\sk{प्रतिपालयति}} (\emph{pratipālayati}) is built on \emph{prati} and the root √pāl, \textquotedblleft{}to guard, to keep\textquotedblright{}, which also gives \emph{pālayati}; MW gives \emph{pratipālayati} as to protect, to keep, to observe, and to wait for (MW).
\end{cousinweb}

\begin{grammarlens}[title={What you've built}]
One more word for saying what happens if.
\end{grammarlens}

\subsection*{Your turn}
\begin{itemize}
\raggedright
  \item \emph{Say it:} \emph{pratipālayati}
  \item \emph{Say it:} \emph{pratipālayati}, once more
  \item \emph{Say it:} \emph{yadi saḥ niyamaṃ pratipālayati, tarhi sarve taṃ viśvasanti.} — If he keeps the rule, everyone trusts him.
  \item \emph{Recall:} say the Sanskrit for turns down, then the Sanskrit for consent, then say \emph{pratipālayati} again
\end{itemize}

\begin{tcolorbox}[breakable,colback=teal!4,colframe=teal!35!black,title={Before you move on}]
What does \sk{प्रतिपालयति} mean? (\textquotedblleft{}Keeps to; abides by\textquotedblright{}.) Say it once more.
\end{tcolorbox}
\section[\texorpdfstring{\sk{सत्यसन्धः} (satyasandhaḥ)}{सत्यसन्धः (satyasandhaḥ)}]{\sk{सत्यसन्धः} (satyasandhaḥ) — true to one's word}

\label{lesson:SA-C535-satyasandhah}



\begin{quote}
Before the new one: say the Sanskrit for consent, then the Sanskrit for keeps to.
\end{quote}

\subsection*{You'll want to know: \sk{सत्यसन्धः}}
\textbf{\sk{सत्यसन्धः}} — \emph{satyasandhaḥ} — \textquotedblleft{}true to one's word\textquotedblright{}.

\begin{quote}
\sk{यदि} \sk{सः} \sk{सत्यसन्धः} \sk{अस्ति}, \sk{तर्हि} \sk{सः} \sk{अवश्यम्} \sk{आगच्छति।}
\end{quote}

\begin{quote}
\emph{yadi saḥ satyasandhaḥ asti, tarhi saḥ avaśyam āgacchati.}
\end{quote}

\begin{quote}
If he keeps his word, he will surely come.
\end{quote}

\begin{cousinweb}
\textbf{\sk{सत्यसन्धः}} (\emph{satyasandhaḥ}) is built from \emph{satya}, \textquotedblleft{}true\textquotedblright{}, and \emph{saṃdhā}, \textquotedblleft{}a compact, a promise\textquotedblright{}; MW gives \emph{satyasaṃdha} as true to engagements, keeping one's promise (MW).
\end{cousinweb}

\begin{grammarlens}[title={What you've built}]
One more word for saying what happens if.
\end{grammarlens}

\subsection*{Your turn}
\begin{itemize}
\raggedright
  \item \emph{Say it:} \emph{satyasandhaḥ}
  \item \emph{Say it:} \emph{satyasandhaḥ}, once more
  \item \emph{Say it:} \emph{yadi saḥ satyasandhaḥ asti, tarhi saḥ avaśyam āgacchati.} — If he keeps his word, he will surely come.
  \item \emph{Recall:} say the Sanskrit for consent, then the Sanskrit for keeps to, then say \emph{satyasandhaḥ} again
\end{itemize}

\begin{tcolorbox}[breakable,colback=teal!4,colframe=teal!35!black,title={Before you move on}]
What does \sk{सत्यसन्धः} mean? (\textquotedblleft{}True to one's word\textquotedblright{}.) Say it once more.
\end{tcolorbox}
\section[\texorpdfstring{\sk{विश्वस्तः} (viśvastaḥ)}{विश्वस्तः (viśvastaḥ)}]{\sk{विश्वस्तः} (viśvastaḥ) — trusted; trustworthy}

\label{lesson:SA-C535-visvastah}



\begin{quote}
Before the new one: say the Sanskrit for keeps to, then the Sanskrit for true to one's word.
\end{quote}

\subsection*{You'll want to know: \sk{विश्वस्तः}}
\textbf{\sk{विश्वस्तः}} — \emph{viśvastaḥ} — \textquotedblleft{}trusted; trustworthy\textquotedblright{}.

\begin{quote}
\sk{यदि} \sk{मित्रं} \sk{विश्वस्तम्} \sk{अस्ति}, \sk{तर्हि} \sk{भयं} \sk{न} \sk{अस्ति।}
\end{quote}

\begin{quote}
\emph{yadi mitraṃ viśvastam asti, tarhi bhayaṃ na asti.}
\end{quote}

\begin{quote}
If a friend is trustworthy, there is nothing to fear.
\end{quote}

\begin{cousinweb}
\textbf{\sk{विश्वस्तः}} (\emph{viśvastaḥ}) is built on \emph{vi} and the root √śvas, \textquotedblleft{}to breathe\textquotedblright{}, as its past participle, the root of \emph{viśvasiti}, \textquotedblleft{}trusts\textquotedblright{}; MW gives \emph{viśvasta} as full of confidence, and as trusted, faithful (MW).
\end{cousinweb}

\begin{grammarlens}[title={What you've built}]
One more word for saying what happens if.
\end{grammarlens}

\subsection*{Your turn}
\begin{itemize}
\raggedright
  \item \emph{Say it:} \emph{viśvastaḥ}
  \item \emph{Say it:} \emph{viśvastaḥ}, once more
  \item \emph{Say it:} \emph{yadi mitraṃ viśvastam asti, tarhi bhayaṃ na asti.} — If a friend is trustworthy, there is nothing to fear.
  \item \emph{Recall:} say the Sanskrit for keeps to, then the Sanskrit for true to one's word, then say \emph{viśvastaḥ} again
\end{itemize}

\begin{tcolorbox}[breakable,colback=teal!4,colframe=teal!35!black,title={Before you move on}]
What does \sk{विश्वस्तः} mean? (\textquotedblleft{}Trusted; trustworthy\textquotedblright{}.) Say it once more.
\end{tcolorbox}
\section[\texorpdfstring{\sk{भंगः} (bhaṅgaḥ)}{भंगः (bhaṅgaḥ)}]{\sk{भंगः} (bhaṅgaḥ) — a breaking}

\label{lesson:SA-C535-bhangah}



\begin{quote}
Before the new one: say the Sanskrit for true to one's word, then the Sanskrit for trusted.
\end{quote}

\subsection*{You'll want to know: \sk{भंगः}}
\textbf{\sk{भंगः}} — \emph{bhaṅgaḥ} — \textquotedblleft{}a breaking\textquotedblright{}.

\begin{quote}
\sk{यदि} \sk{नियमस्य} \sk{भंगः} \sk{भवति}, \sk{तर्हि} \sk{अध्यापकः} \sk{कुपितः} \sk{भवति।}
\end{quote}

\begin{quote}
\emph{yadi niyamasya bhaṅgaḥ bhavati, tarhi adhyāpakaḥ kupitaḥ bhavati.}
\end{quote}

\begin{quote}
If a rule is broken, the teacher gets angry.
\end{quote}

A masculine noun.

\begin{cousinweb}
\textbf{\sk{भंगः}} (\emph{bhaṅgaḥ}) comes from the root √bhañj, \textquotedblleft{}to break\textquotedblright{}; MW gives \emph{bhaṅga} as breaking, splitting, breaking down (MW).
\end{cousinweb}

\begin{grammarlens}[title={What you've built}]
One more word for saying what happens if.
\end{grammarlens}

\subsection*{Your turn}
\begin{itemize}
\raggedright
  \item \emph{Say it:} \emph{bhaṅgaḥ}
  \item \emph{Say it:} \emph{bhaṅgaḥ}, once more
  \item \emph{Say it:} \emph{yadi niyamasya bhaṅgaḥ bhavati, tarhi adhyāpakaḥ kupitaḥ bhavati.} — If a rule is broken, the teacher gets angry.
  \item \emph{Recall:} say the Sanskrit for true to one's word, then the Sanskrit for trusted, then say \emph{bhaṅgaḥ} again
\end{itemize}

\begin{tcolorbox}[breakable,colback=teal!4,colframe=teal!35!black,title={Before you move on}]
What does \sk{भंगः} mean? (\textquotedblleft{}A breaking\textquotedblright{}.) Say it once more.
\end{tcolorbox}
\section[\texorpdfstring{\sk{वचनम्} (vacanam)}{वचनम् (vacanam)}]{\sk{वचनम्} (vacanam) — one's word; a statement}

\label{lesson:SA-C535-vacanam}



\begin{quote}
Before the new one: say the Sanskrit for trusted, then the Sanskrit for a breaking.
\end{quote}

\subsection*{You'll want to know: \sk{वचनम्}}
\textbf{\sk{वचनम्}} — \emph{vacanam} — \textquotedblleft{}one's word; a statement\textquotedblright{}.

\begin{quote}
\sk{यदि} \sk{भवान्} \sk{वचनं} \sk{ददाति}, \sk{तर्हि} \sk{तत्} \sk{प्रतिपालयतु।}
\end{quote}

\begin{quote}
\emph{yadi bhavān vacanaṃ dadāti, tarhi tat pratipālayatu.}
\end{quote}

\begin{quote}
If you give your word, keep to it.
\end{quote}

A neuter noun.

\begin{cousinweb}
\textbf{\sk{वचनम्}} (\emph{vacanam}) comes from the root √vac, \textquotedblleft{}to speak\textquotedblright{}, with the ending \emph{-ana}; MW gives \emph{vacana} as an utterance, a statement, a word (MW).
\end{cousinweb}

\begin{grammarlens}[title={What you've built}]
One more word for saying what happens if.
\end{grammarlens}

\subsection*{Your turn}
\begin{itemize}
\raggedright
  \item \emph{Say it:} \emph{vacanam}
  \item \emph{Say it:} \emph{vacanam}, once more
  \item \emph{Say it:} \emph{yadi bhavān vacanaṃ dadāti, tarhi tat pratipālayatu.} — If you give your word, keep to it.
  \item \emph{Recall:} say the Sanskrit for trusted, then the Sanskrit for a breaking, then say \emph{vacanam} again
\end{itemize}

\begin{tcolorbox}[breakable,colback=teal!4,colframe=teal!35!black,title={Before you move on}]
What does \sk{वचनम्} mean? (\textquotedblleft{}One's word; a statement\textquotedblright{}.) Say it once more.
\end{tcolorbox}
